% !TeX program = pdflatex
\documentclass[12pt,a4paper]{article}

\newcommand{\notelanguage}{finnish}
% Shared preamble for course notes and exercise sets.

\usepackage[T1]{fontenc}
\usepackage{lmodern}
\usepackage[english]{babel}

\usepackage[a4paper,margin=30mm]{geometry}
\usepackage{microtype}
\usepackage{graphicx}
\usepackage{enumitem}

\usepackage{amsmath}
\usepackage{amssymb}
\usepackage{amsthm}
\usepackage{mathtools}
\usepackage{aliascnt}

\usepackage[hidelinks,pdfusetitle]{hyperref}

\numberwithin{equation}{section}
\setlist{itemsep=0.25em,topsep=0.5em}

\theoremstyle{plain}
\newtheorem{theorem}{Theorem}[section]
\newaliascnt{lemma}{theorem}
\newtheorem{lemma}[lemma]{Lemma}
\aliascntresetthe{lemma}
\newaliascnt{proposition}{theorem}
\newtheorem{proposition}[proposition]{Proposition}
\aliascntresetthe{proposition}
\newaliascnt{corollary}{theorem}
\newtheorem{corollary}[corollary]{Corollary}
\aliascntresetthe{corollary}

\theoremstyle{definition}
\newaliascnt{definition}{theorem}
\newtheorem{definition}[definition]{Definition}
\aliascntresetthe{definition}
\newaliascnt{example}{theorem}
\newtheorem{example}[example]{Example}
\aliascntresetthe{example}
\newaliascnt{problem}{theorem}
\newtheorem{problem}[problem]{Problem}
\aliascntresetthe{problem}

\theoremstyle{remark}
\newaliascnt{remark}{theorem}
\newtheorem{remark}[remark]{Remark}
\aliascntresetthe{remark}

\usepackage[nameinlink,noabbrev]{cleveref}

\crefname{theorem}{theorem}{theorems}
\crefname{lemma}{lemma}{lemmas}
\crefname{proposition}{proposition}{propositions}
\crefname{corollary}{corollary}{corollaries}
\crefname{definition}{definition}{definitions}
\crefname{example}{example}{examples}
\crefname{problem}{problem}{problems}
\crefname{remark}{remark}{remarks}

\DeclareMathOperator{\im}{im}
\DeclareMathOperator{\rank}{rank}
\DeclarePairedDelimiter{\abs}{\lvert}{\rvert}
\DeclarePairedDelimiter{\norm}{\lVert}{\rVert}
\DeclarePairedDelimiter{\set}{\lbrace}{\rbrace}


\usepackage{tikz}

\title{Ryhmätyö 1 ratkaisut}
\author{}
\date{}

\begin{document}

\maketitle

\noindent

\section{Ratkaisut}

\begin{exercise}
  \leavevmode\par
\begin{enumerate}[label=\alph*), nosep]
  \item Mitkä luvut toteuttavat epäyhtälön \(\abs{2x - 4} < 7\)? Piirrä
    kuva tilanteesta. Kirjoita vastaus muodossa ``\(x \in (a, b)\)`` ja
    muodossa ``\(c < x < d\)``.
  \item Miten vastaus muuttuu, jos epäyhtälössä \(\abs{2x - 4} < 7\) merkki
    \(<\) korvataan merkillä \(\leq\)?
  \item Voiko epäyhtälöä muokata niin, että vastaukseksi saadaan
    väli \((a, b]\) tai \([a, b)\)?
  \item Onko olemassa toista epäyhtälöä, joka antaa saman välin kuin
    tapauksessa (a)?
\end{enumerate}
\end{exercise}

\begin{solution}
\leavevmode\par
\begin{enumerate}[label=\alph*), nosep]
\item Sieventämällä ja itseisarvon määritelmää hyödyntäen saadaan
  \[-7 < 2x - 4 < 7 \iff -\frac{3}{2} < x < \frac{11}{2},\]
  eli \(x \in \left(-\frac{3}{2}, \frac{11}{2}\right)\).
  \begin{center}
    \begin{tikzpicture}[x=1.2cm]
      \draw[<->] (-3, 0) -- (7, 0) node[right] {$x$};
      \draw[very thick] (-1.5, 0) -- (5.5, 0);
      \draw[fill=white, thick] (-1.5, 0) circle (2.5pt);
      \draw[fill=white, thick] (5.5, 0) circle (2.5pt);
      \draw (-1.5, 0.12) -- (-1.5, -0.12)
      node[below=4pt] {$-\frac{3}{2}$};
      \draw (2, 0.12) -- (2, -0.12)
      node[below=4pt] {$2$};
      \draw (5.5, 0.12) -- (5.5, -0.12)
      node[below=4pt] {$\frac{11}{2}$};
    \end{tikzpicture}
  \end{center}
\item Vaihtamalla epäyhtälön merkiksi \(\leq\), saadaan luvulle \(x\) ehdot
  \[x \in \left[-\frac{3}{2}, \frac{11}{2}\right] \iff -\frac{3}{2} \leq x \leq
  \frac{11}{2}.\]
  \begin{center}
    \begin{tikzpicture}[x=1.2cm]
      \draw[<->] (-3, 0) -- (7, 0) node[right] {$x$};
      \draw[very thick] (-1.5, 0) -- (5.5, 0);
      \draw[fill=white, thick] (-1.5, 0) circle (2.5pt);
      \fill (-1.5, 0) circle (2.5pt);
      \draw[fill=white, thick] (5.5, 0) circle (2.5pt);
      \fill (5.5, 0) circle (2.5pt);
      \draw (-1.5, 0.12) -- (-1.5, -0.12)
      node[below=4pt] {$-\frac{3}{2}$};
      \draw (2, 0.12) -- (2, -0.12)
      node[below=4pt] {$2$};
      \draw (5.5, 0.12) -- (5.5, -0.12)
      node[below=4pt] {$\frac{11}{2}$};
    \end{tikzpicture}
  \end{center}
\item Kyllä. Puoliavoin väli saadaan käyttämällä eri päätepisteissä eri
  epäyhtälömerkkejä. Esimerkiksi
  \[-7 < 2x - 4 \leq 7\]
  antaa
  \[x \in \left(-\frac{3}{2}, \frac{11}{2}\right].\]
  Vastaavasti
  \[-7 \leq 2x - 4 < 7\]
  antaa
  \[x \in \left[-\frac{3}{2}, \frac{11}{2}\right).\]
  Pelkkä muoto \(\abs{2x - 4} < 7\) tai \(\abs{2x - 4} \leq 7\) ei voi antaa
  puoliavointa väliä, koska itseisarvo käsittelee molempia rajoja symmetrisesti.
\item Kyllä, esimerkiksi epäyhtälö
  \[\abs{x - 2} < \frac{7}{2}.\]
  Tämän voi nähdä ekuivalensseista
  \begin{align*}
    \abs{x - 2} < \frac{7}{2} &\iff 2\abs{x - 2} < 7 \\
    &\iff \abs{2x - 4} < 7.
  \end{align*}
  Joten epäyhtälöilla \(\abs{2x - 4} < 7\) ja \(\abs{x - 2} < \frac{7}{2}\)
  on identtinen ratkaisujoukko, eli luvut \(x \in \left(-\frac{3}{2},
  \frac{11}{2}\right)\).
\end{enumerate}
\end{solution}

\begin{exercise}
Määritellään joukko \(E\) asettamalla
\[E = \left\{m \in \RR : \abs{x^2 - 1} \leq m\abs{x - 1} \quad
\text{kaikilla} \quad x \in (-1, 3) \right\}.\]
\begin{enumerate}[label=\alph*), nosep]
\item Mitä eri joukon \(E\) määrittelevän lausekkeen osat tarkoittavat.
(Mikä on esimerkiksi kaksoispisteen merkitys? Mitä muita osia lausekkeessa on?)
\item Onko joukossa \(E\) suurinta tai pienintä alkiota?
\item Miten tehtävien (a) ja (b) vastaukset muuttuvat, jos ehto ``\(x
\in (-1, 3)\)`` korvataan ehdolla ``\(x \in [-1, 3]\)``?
\end{enumerate}
\end{exercise}

\begin{solution}
\leavevmode\par
\begin{enumerate}[label=\alph*), nosep]
\item Joukon \(E\) määrittelevä lause käyttää yleistä tapaa määrittää joukko.
Tarkastellaan ensin yleistä tapausta. Olkoon
\begin{equation}\label{set:builder}
S = \left\{x \in A : P(x)\right\}.
\end{equation}
Määritelmässä \ref{set:builder} \(S\) on määritellyn joukon \emph{nimi}.
Merkinnän osa \(x \in A\) kertoo, että joukko \(S\) koostuu joukon \(A\)
alkioista. Merkintä \(:\) voidaan lukea ``siten, että``, ``jolle``
tai ``joilla``. Merkintä \(P(x)\) kuvaa \emph{ominaisuutta},
jonka perusteella joukosta \(A\) valitaan ne alkiot, jotka kuuluvat joukkoon
\(S\). Kokonaisuudessa siis, joukko \(S\) koostuu joukon \(A\) alkioista \(x\),
jotka toteuttavat ominaisuuden \(P(x)\).

Voimme tämän yleisen kuvauksen perusteella tutkia joukon \(E\) määritelmää.
Yllä olevan perusteella, joukko \(E\) koostuu luvuista \(m \in \RR\),
jotka toteuttavat seuraavan ehdon:
\[\abs{x^2 - 1} \leq m\abs{x - 1} \quad \text{kaikilla} \quad x \in (-1, 3).\]
Tässä \(m\) on joukon \(E\) alkioehdokas, kun taas \(x\) käy läpi kaikki välin
\((-1, 3)\) luvut. Ehto tarkoittaa, että epäyhtälön
\[\abs{x^2 - 1} \leq m\abs{x - 1}\]
on oltava voimassa jokaisella \(x \in (-1, 3)\). Joukkoon \(E\) kuuluu
täsmälleen ne reaaliluvut \(m\), joilla tämä toteutuu.
\item Sieventämällä epäyhtälöä saadaan
\[\abs{x^2 - 1} \leq m\abs{x - 1} \iff \abs{x + 1}\abs{x - 1} \leq
m\abs{x - 1},\]
joten kun \(x \neq 1\), saadaan
\[\abs{x + 1} \leq m.\]
On voimassa \(x \in (-1, 3)\), joten \(\abs{x + 1} \in (0, 4).\) Kun taas
\(x = 1\), epäyhtälö saa muodon \(0 \leq 0\), joten tämä tapaus ei anna
lisäehtoja luvulle \(m\).
Koska epäyhtälön on oltava voimassa kaikilla \(x \in (-1,3)\) ja
\(\abs{x + 1}\) saa kaikki arvot väliltä \((0, 4)\), luvun \(m\) on oltava
välin \((0,4)\) yläraja. Siis \(m \ge 4\), joten
\[E=\{m \in \RR : m \geq 4\} = [4, \infty).\]
Joukon \(E\) pienin alkio on \(4\), eikä sillä ole suurinta alkiota.
\item On siis määritelty
\[E = \left\{m \in \RR : \abs{x^2 - 1} \leq m\abs{x - 1} \quad
\text{kaikilla} \quad x \in [-1, 3] \right\}.\]
Samalla logiikalla kuin kohdassa (b), saamme
\[\abs{x + 1} \leq m, \quad x \in [-1, 3] \setminus{\{1\}}.\]
Koska \(\abs{x + 1} \leq 4\) kaikilla \(x \in [-1, 3]\) ja \(\abs{x + 1} = 4\),
kun \(x = 3\), on oltava \(m \geq 4\). Toisaalta jokainen
\(m \geq 4\) toteuttaa ehdon, sillä \(\abs{x + 1} \leq 4 \leq m\)
kaikilla \(x \in [-1, 3] \setminus{\{1\}}\), ja \(x = 1\) toteuttaa
epäyhtälön triviaalisti. Täten
\[E = \{m \in \RR : m \geq 4\} = [4, \infty).\]
Täten joukon \(E\) pienin alkio on jälleen \(4\), eikä suurinta alkiota ole.
Merkittävin ero kohtaan (b) on, että nyt arvo \(4\) saavutetaan:
\(\abs{x + 1} = 4\), kun \(x = 3\). Kohdassa (b), arvo \(4\) oli vain
joukon \((0, 4)\) \emph{supremum}.
\end{enumerate}
\end{solution}

\end{document}
